% =====================================================================
\section{Setup: Two Directed Distances and a Decomposition}
\label{sec:setup}

\begin{figure}[t]
\centering
\includegraphics[width=\linewidth]{figures/splats/composite_rgb_motion.png}
\caption{\textbf{Motion fingerprints (BFV splats) and RGB samples.}
\emph{Top block:} each panel aggregates a dataset's joint flow vectors
$\mathbf{f}=[\hat x,\hat y,\Delta\hat x,\Delta\hat y]$, clusters them spatially
(non-uniform, data-driven clusters), and renders each cluster as a directional
splat: color encodes dominant flow direction (legend), splat extent encodes the
cluster's spread; training sets (upper row) above benchmarks (lower row). The
first three columns hold the same dataset in both rows (train/eval); their
fingerprints match, confirming the signature is a property of the data rather
than of a particular split. Training sets and benchmarks otherwise occupy
visibly different regions of motion space. \emph{Bottom block:} one RGB sample
per dataset, same column order, for qualitative reference. The two directed
distances defined below quantify the two distinct ways motion supports can fail
to align.}
\label{fig:splats}
\end{figure}

We represent each dataset by the distribution of its motion and ask how well a
training set's distribution covers a target's. This section defines that
representation, the two directed distances that are our central measurement,
and the decomposition we use to turn them into predictions.

\noindent\textbf{Transfer study.}
We study four architecture families: the semantic matcher CATs++
\cite{cho2022cats++}, the optical-flow networks RAFT \cite{teed2020raft} and
FlowFormer \cite{huang2022flowformer}, and the dense-matching network GLU-Net
\cite{truong2020glunet}. For CATs++, GLU-Net, and FlowFormer the image-feature
backbone (ImageNet ResNet-101, VGG, and a pretrained Twins-SVT transformer,
respectively) is either initialized from those features or trained from
scratch, and either frozen or trained jointly with the correspondence head,
giving four configurations each; RAFT is trained from scratch by construction
and is one. The resulting \textbf{13 backbone configurations} are each trained
on 11 pure synthetic sources (\textbf{142 trained models}) and evaluated by peak \pck@0.05 on 9 benchmarks spanning
optical flow, dense, and semantic matching (KITTI-2012/2015 \cite{kitti_flow},
FlyingThings \cite{mayer2016FlyingThingsDataset}, PointOdyssey
\cite{pointodyssey}, TSS, PF-Pascal/PF-Willow \cite{pfpascal}, SPair-71k
\cite{spair71k}, and a held-out split of our procedural source), for
\textbf{1{,}278 source--benchmark--configuration transfer measurements}. Six of
the nine targets are real images (KITTI-2012/2015 and the four semantic sets);
FlyingThings, PointOdyssey, and the procedural split are synthetic targets,
included for their controlled motion and exact labels (so real-image \emph{motion}
transfer is carried mainly by KITTI, a scope point we return to in
Sec.~\ref{sec:limits}). The
coverage score itself depends only on the (source, target-benchmark) pair, but a
source's realized transfer also depends on the model, so we evaluate every ranking
claim within a fixed (benchmark, architecture, configuration) \emph{context}; this
same context indexes the level term of the absolute predictor (Sec.~\ref{sec:absolute}).
We report \emph{peak} \pck{} (the best checkpoint per run) so that data quality is
not confounded with convergence speed: a fixed-budget last-checkpoint number would
penalize sources that need a different schedule. Because peak selection is applied
identically to every source in a context, and the selection rule only ever ranks
sources \emph{within} a context, it is a constant, monotone treatment that cannot
change which source ranks higher---it leaves the ranking $\rho$ untouched and can
inflate only the absolute level, which we predict separately and openly attribute
in part to a realism effect (Sec.~\ref{sec:absolute}). Throughout we stratify the nine
benchmarks into \emph{real-motion} targets (KITTI-2012/2015, FlyingThings,
PointOdyssey, and the procedural split), whose ground truth is dense pixel
motion, and \emph{semantic} targets (SPair-71k, PF-Pascal, PF-Willow, TSS),
whose ground truth is sparse keypoint correspondence across object instances;
we report coverage stratified by both, as it behaves differently on the two
target types (Sec.~\ref{sec:law}).

\noindent\textbf{Each architecture in its designed regime.} CATs++, GLU-Net, and
FlowFormer are designed to run on a \emph{pretrained} image backbone; RAFT is the
one architecture built to train end-to-end from scratch with no pretrained
backbone. We accordingly analyze each in its intended regime: the three
backbone-dependent matchers with their pretrained backbones (encoder trained or
frozen), and RAFT from scratch. Training the backbone-dependent matchers from
scratch is off-design: their randomly-initialized backbones mode-collapse to a
near-static flow prior that ignores the image pair, so a \emph{data}-selection
signal cannot exist there at all---a network that ignores its input cannot reveal
which source trained it. We confirm this with two motion-free checks: a
source-sensitivity probe (a collapsed model barely changes its prediction when the
source image is swapped) and the test of whether a from-scratch encoder beats its
own random initialization ($\Delta_{\text{enc}}{>}0$). The off-design
configurations fail both, and we report them in full in the supplement. Because
this partition follows from architecture design and is fixed before any motion
distance is measured, it introduces no outcome-based selection.

\noindent\textbf{Motion representation.}
We summarize each dataset as a point cloud of flow vectors in a normalized
four-dimensional space of image position and displacement,
\begin{equation}
\mathbf{f}_{\text{flow}} = [\hat{x},\, \hat{y},\, \Delta\hat{x},\,
\Delta\hat{y}] \in [-1,1]^4,
\label{eq:bfv}
\end{equation}
following the bag-of-flow-vectors representation (Fig.~\ref{fig:splats}). This
retains \emph{where} motion occurs and \emph{how} content moves while remaining
invariant to image resolution and free of any spatial binning. For the
appearance control, we build the analogous representation from DINOv3 patch
features, reduced by PCA and $\ell_2$-normalized \cite{simeoni2025dinov3}.

\noindent\textbf{Two directed distances.}
The match between a training cloud $A$ and a target cloud $B$ is directional,
and we measure each direction by the mean nearest-neighbor distance
\begin{equation}
d_{A \to B} \;=\; \frac{1}{|A|}\sum_{a \in A}\; \min_{b \in B}\,
\lVert a - b \rVert_2^2 .
\label{eq:dirdist}
\end{equation}
Writing $\dtb$ for the train-to-target direction and $\dbt$ for the reverse,
$\dtb$ is large when much of the training distribution lies far from anything
the target uses (\textbf{off-target mass}), and $\dbt$ is large when target
motions have no nearby training example (\textbf{missing support}).

Each direction can be estimated in several ways. Our default,
Eq.~(\ref{eq:dirdist}), is a \emph{support-geometry} estimator: it depends only
on how close the two point sets lie, not on their local densities. We compare it
(supp.) against two alternatives: $\epsilon$-coverage, the fraction of one cloud
within a radius of the other (also support-based), and a $k$-nearest-neighbor
estimate of the KL divergence (density-based). We prefer the support-based
estimators for procedurally generated data on empirical grounds: the
nearest-neighbor and $\epsilon$-coverage estimates converge as the sample grows,
whereas the KL estimate fails to stabilize at any budget we tried (supp.). We
suspect this is because rendering more frames of the same motion changes a
density estimate while leaving the support unchanged, but we do not test that
mechanism directly; the instability alone makes the density estimator unreliable
here. Of the three, nearest-neighbor is the most informative predictor and is our
default, $\epsilon$-coverage is close behind, and KL is the least stable and the
weakest; that all three nonetheless agree on the central finding
(Sec.~\ref{sec:law}) shows it reflects the \emph{direction} of the comparison
rather than the estimator. As direction-blind baselines we also use the symmetric
average of the two directions and two distribution distances, FID
\cite{heusel2018ganstrainedtimescaleupdate} and sliced Wasserstein-2.

\noindent\textbf{Correlation convention.} We report Spearman $\rho$ for every
\emph{ranking} claim (which source is better---the quantity the selection rule
acts on) and Pearson $r$ (with MAE) only for \emph{absolute-\pck} prediction
(Sec.~\ref{sec:absolute}). The pooled scatter (Fig.~\ref{fig:coverage_law_a}) shows
Pearson $r$ on within-context-standardized values for visualization; the
load-bearing ranking statistic is the Spearman $\rho$ of Table~\ref{tab:law}.

